\section{Introduction}
\label{sec:intro}



The Multi-modal Aerial View Image Challenge for cross domain image translation~\cite{lowMultimodalAerialView2024,lowMultimodalAerialView2023,bowald3rdMultimodalAerial2025} is a leading challenge held in conjunction with the Perception Beyond the Visible Spectrum (PBVS) workshop from 2023 to 2025. Data collected across multiple sensing modalities---such as Electro-Optical (EO), Infrared (IR), and Synthetic Aperture Radar (SAR)---each capture distinct characteristics and measurements of objects, providing complementary perspectives on the same scene. However, data coverage limitations and modality-specific constraints often hinder the full utilization of these diverse sources. For instance, while SAR sensors offer all-weather capabilities and can penetrate atmospheric obstructions, their imagery is complex to interpret and publicly available datasets are scarce. Conversely, widely available EO imagery is limited by lighting conditions and atmospheric visibility. Domain translation seeks to bridge this gap by converting one modality into another, enabling cross-modal data synthesis and enhancing sensor interoperability. By focusing on translation tasks such as RGB$\rightarrow$IR, SAR$\rightarrow$EO, SAR$\rightarrow$IR, and SAR$\rightarrow$RGB, the MAVIC-T challenge aims to bolster existing data volume, increase data diversity across regions, and serve as a platform for advancing conditioned image synthesis techniques, ultimately fostering the development of more adaptable and robust models for multi-modal image analysis.

% paragraph: talk about the difficulty to do image translation across above different modality (with proper work citation)
interpret SAR need special knowledge~\cite{majumder2020deep}.

% paragraph: our solution
In this paper, we propose [YourMethodName], a novel [Diffusion/GAN/Transformer] framework designed to [solve specific problem]. Our approach incorporates a [Specific Module, e.g., 'Latent Consistency Constraint'] and a [Specific Loss] to ensure that the translated images maintain the precise geometric layout of the source modality while achieving high photorealism.

% (optional) paragraph: Key Findings & Results
% Our findings indicate that [Specific Technique] significantly reduces the FID score by [X%]. Qualitatively, our model demonstrates a superior ability to reconstruct urban structures in SAR-to-RGB tasks compared to the baseline models provided in the challenge. Furthermore, our solution placed [Rank] in the [Task Name] leaderboard, proving its robustness across diverse geographical regions.

% paragraph: Contributions